\documentclass[aspectratio=169]{beamer}

\usepackage{fontspec}
\usepackage{amsmath}
\usepackage{xcolor}

\usetheme{Madrid}
\setmainfont{DejaVu Sans}
\setsansfont{DejaVu Sans}

\title{Chuẩn hóa kích thước ảnh MRI trong PKA}
\subtitle{Ý tưởng lý thuyết của resize và resampling}
\author{P2VG preprocessing}
\date{\today}

\begin{document}

\begin{frame}
  \titlepage
\end{frame}

\begin{frame}{Vấn đề ban đầu}
  Ảnh MRI từ nhiều bệnh nhân thường không đồng nhất:

  \begin{itemize}
    \item Số lát cắt khác nhau.
    \item Kích thước mặt phẳng ảnh khác nhau.
    \item Khoảng cách voxel khác nhau.
    \item Hướng trục ảnh có thể không hoàn toàn giống nhau.
  \end{itemize}

  \vspace{0.3cm}
  Nếu đưa trực tiếp vào model, các ảnh sẽ không có cùng hình dạng tensor và không cùng ý nghĩa không gian.
\end{frame}

\begin{frame}{Mục tiêu chuẩn hóa}
  Mục tiêu là đưa mọi ảnh về một biểu diễn thống nhất:

  \[
    256 \times 256 \times 32
  \]

  \begin{itemize}
    \item Hai chiều đầu là mặt phẳng ảnh.
    \item Chiều cuối là số lát cắt theo chiều sâu.
    \item Tất cả ảnh có cùng kích thước tensor để đưa vào mạng học sâu.
  \end{itemize}
\end{frame}

\begin{frame}{Resize không chỉ là đổi số pixel}
  Với ảnh y khoa 3D, mỗi voxel đại diện cho một thể tích thật trong cơ thể.

  \begin{itemize}
    \item Resize thông thường chỉ quan tâm đến kích thước mảng.
    \item Resampling y khoa cần quan tâm thêm spacing vật lý.
    \item Spacing cho biết mỗi voxel tương ứng bao nhiêu mm ngoài đời thật.
  \end{itemize}

  \vspace{0.3cm}
  Vì vậy pipeline tách thành hai ý tưởng: chuẩn hóa spacing trước, rồi chuẩn hóa shape sau.
\end{frame}

\begin{frame}{Bước 1: Chuẩn hóa spacing}
  Ảnh đầu vào có spacing ban đầu:

  \[
    (s_x, s_y, s_z)
  \]

  Spacing mục tiêu mặc định:

  \[
    (1.0,\ 1.0,\ 4.0)\ \text{mm}
  \]

  \begin{itemize}
    \item Hai chiều trong mặt phẳng ảnh có độ phân giải 1 mm.
    \item Chiều sâu có spacing 4 mm, phù hợp với dữ liệu lát cắt dày hơn.
    \item Bước này giúp ảnh của các bệnh nhân có tỷ lệ không gian nhất quán hơn.
  \end{itemize}
\end{frame}

\begin{frame}{Ý nghĩa của resampling}
  Resampling là quá trình lấy mẫu lại ảnh trên một lưới voxel mới.

  \begin{itemize}
    \item Nếu spacing mới nhỏ hơn, số voxel có thể tăng lên.
    \item Nếu spacing mới lớn hơn, số voxel có thể giảm đi.
    \item Giá trị voxel mới được ước lượng từ các voxel lân cận.
  \end{itemize}

  \vspace{0.3cm}
  Với ảnh MRI intensity, nội suy tuyến tính giúp giữ chuyển tiếp cường độ tương đối mượt.
\end{frame}

\begin{frame}{Bước 2: Chuẩn hóa số lát cắt}
  Sau khi spacing đã được chuẩn hóa, số lát cắt vẫn có thể khác nhau.

  \[
    D \rightarrow 32
  \]

  \begin{itemize}
    \item Nếu ảnh có nhiều hơn 32 lát, ảnh được nén theo chiều sâu.
    \item Nếu ảnh có ít hơn 32 lát, ảnh được nội suy thêm lát trung gian.
    \item Mục tiêu là mọi ảnh có cùng chiều depth để model xử lý ổn định.
  \end{itemize}
\end{frame}

\begin{frame}{Bước 3: Chuẩn hóa mặt phẳng ảnh}
  Hai chiều trong mặt phẳng được đưa về:

  \[
    H \times W \rightarrow 256 \times 256
  \]

  \begin{itemize}
    \item Giúp mỗi lát cắt có cùng kích thước.
    \item Giảm khác biệt giữa scanner, protocol và kích thước crop ban đầu.
    \item Tạo tensor đầu vào cố định cho ViT 3D.
  \end{itemize}
\end{frame}

\begin{frame}{Vì sao chọn \(256 \times 256 \times 32\)?}
  Kích thước này là một thỏa hiệp giữa thông tin ảnh và chi phí tính toán.

  \begin{itemize}
    \item \(256 \times 256\): đủ chi tiết cho cấu trúc cột sống thắt lưng.
    \item \(32\) lát: giữ thông tin 3D nhưng không làm bộ nhớ GPU quá lớn.
    \item Kích thước cố định giúp batch training đơn giản và ổn định hơn.
  \end{itemize}
\end{frame}

\begin{frame}{Giữ thứ tự trục ảnh}
  Pipeline hiện tại dùng ý tưởng \textit{preserve axis}.

  \begin{itemize}
    \item Không ép ảnh về một orientation thế giới mới.
    \item Giữ thứ tự trục mảng như dữ liệu đầu vào sau khi resize.
    \item Phù hợp khi toàn bộ dataset đã có quy ước trục tương đối ổn định.
  \end{itemize}

  \vspace{0.3cm}
  Điều này giúp tránh thay đổi bất ngờ vị trí tương đối của sagittal, axial và fused image.
\end{frame}

\begin{frame}{Vai trò của affine}
  Affine mô tả quan hệ giữa tọa độ voxel và tọa độ vật lý.

  \begin{itemize}
    \item Khi resize ảnh, số voxel thay đổi.
    \item Khi số voxel thay đổi, spacing hiệu dụng cũng thay đổi.
    \item Do đó affine cần được cập nhật để metadata không bị sai.
  \end{itemize}

  \vspace{0.3cm}
  Nói ngắn gọn: dữ liệu ảnh thay đổi thì thông tin không gian đi kèm cũng phải thay đổi.
\end{frame}

\begin{frame}{Tại sao cần cùng kích thước cho model?}
  Mạng học sâu thường cần tensor đầu vào có shape cố định trong cùng batch.

  \begin{itemize}
    \item ViT 3D chia ảnh thành các patch 3D.
    \item Kích thước ảnh cố định giúp số patch ổn định.
    \item Số patch ổn định giúp projector và LLM nhận chuỗi image token nhất quán.
  \end{itemize}

  \vspace{0.3cm}
  Vì vậy resize là bước cầu nối giữa ảnh MRI thô và mô hình vision-language.
\end{frame}

\begin{frame}{Tóm tắt}
  Cơ chế thay đổi kích thước ảnh gồm ba ý chính:

  \begin{enumerate}
    \item Chuẩn hóa spacing để các voxel có ý nghĩa vật lý gần giống nhau.
    \item Chuẩn hóa shape về \(256 \times 256 \times 32\) để model nhận input cố định.
    \item Cập nhật thông tin không gian để file ảnh sau xử lý vẫn nhất quán.
  \end{enumerate}

  \vspace{0.3cm}
  Đây không chỉ là resize ảnh, mà là chuẩn hóa biểu diễn không gian của dữ liệu MRI 3D.
\end{frame}

\end{document}
